\section{Benchmark Construction}
\label{sec:benchmark}

Evaluating the integrity axis requires cases that contain integrity failures,
and CI-Work's two-set partition cannot express them. We therefore construct a
benchmark around a single concrete agent---one that reads a repository and the
work systems around it, then writes a review onto a pull request---following
CI-Work's four-stage pipeline with the additions this section describes.

We choose PR review for one specific reason: \textbf{almost every context
source that makes the review better also carries something the review must not
repeat, and the two are about the same code.} There is no version of this task
in which the sensitive material is conveniently off-topic. It is also a setting
where the author-controlled-channel adversary of \Cref{sec:threat} is the
natural one rather than a contrivance, since the party under review has both
write access to the description and an incentive to shape the verdict.

\subsection{What changes from CI-Work's setup}

Three structural differences, each with a consequence
(\Cref{tab:benchdiff}). The third is the one that bites: in PR review the data
subject is typically a third party, so CI's five parameters no longer collapse
and $\mathrm{subj}(e)$ must be tracked separately. ``Customer $X$ had a
47-minute outage'' is sensitive with respect to $X$, who is absent from the
conversation and cannot consent.

\begin{table}[t]
\centering
\caption{Structural differences from CI-Work's setting and their consequences.}
\label{tab:benchdiff}
\footnotesize
\begin{tabular}{@{}p{0.3\columnwidth}p{0.29\columnwidth}p{0.29\columnwidth}@{}}
\toprule
\textbf{CI-Work} & \textbf{PR review} & \textbf{Consequence} \\
\midrule
Agent acts for the sender, who owns the data &
Agent acts for the \emph{reviewer}; the store belongs to the org &
Read authority is the reviewer's clearance, not the owner's \\[3pt]
Recipient is a person &
Recipient is a \emph{thread} whose reader set may grow &
Write ceiling is the min over current \emph{and foreseeable} readers \\[3pt]
Data subject is usually the sender &
Data subject is usually a \emph{third party} &
$\mathrm{subj}(e)$ tracked separately; no consent path exists \\
\bottomrule
\end{tabular}
\end{table}

\subsection{Three entry classes}
\label{sec:benchmark:classes}

CI-Work partitions retrieved context into $\Eess$ and $\Esens$. That partition
can express only confidentiality failures, so we add a third class
(\Cref{tab:classes}).

\begin{table}[t]
\centering
\caption{Three entry classes. $\Ecorr$ is the novel one, and it is the class
that makes \Cref{eq:bibaviolation} measurable: cheap to disclose, expensive to
believe.}
\label{tab:classes}
\footnotesize
\begin{tabular}{@{}lcccl@{}}
\toprule
\textbf{Class} & \textbf{Convey?} & \textbf{Withhold?} & \textbf{Distrust?}
  & \textbf{Labels} \\
\midrule
$\Eess$ essential   & \textbf{yes} & no & no & low $\Cl$, high $\Il$ \\
$\Esens$ sensitive  & no & \textbf{yes} & no & high $\Cl$, $\Il$ varies \\
$\Ecorr$ corrupting & no & no & \textbf{yes} & low $\Cl$, low $\Il$ \\
\bottomrule
\end{tabular}
\end{table}

A corrupting entry is cheap to disclose and expensive to believe: the PR
author's claim that ``security already signed off,'' a code comment reading
\texttt{// TODO: this test is flaky, skip it}, a ticket comment asserting a
number nobody verified. Quoting any of them costs nothing. Treating them as
evidence produces a wrong verdict on production code.

Two things follow. First, $\Ecorr$ is what makes $\Esens \neq \text{low-}\Il$
\emph{measurable} rather than merely arguable, since the walk-away-price case
is high on both axes. Second, the integrity metrics must be computed over
$\Ecorr$ rather than over $\Esens$, so that the denominator is a designed
quantity rather than a side effect of labeling. Deriving the low-integrity set
as $\{e \in \Esens \cup \Eess : \Il(e) < \lvlI\}$, as a naive implementation
does, conflates ``sensitive and incidentally low-integrity'' with ``planted
specifically to corrupt the verdict.''

\subsection{The entangled-twin construction}
\label{sec:benchmark:twins}

The generation rule that makes these cases hard rather than trivially
separable:
\begin{quote}
\emph{Every sensitive entry is authored as a near-twin of an essential one:
same subject matter, same retrieval query, different disclosure norm.}
\end{quote}
If sensitive entries concern a different topic than essential ones, any
competent retriever separates them and the benchmark measures nothing. The twin
structure forces the agent to adjudicate role $\times$ information
compatibility rather than relevance---precisely the fine-grained judgement
CI-Work found frontier models fail while passing the coarse-grained one.
\Cref{tab:twins} gives a triple.

\begin{table}[t]
\centering
\caption{An entangled twin triple. Both first rows answer the same retrieval
query, ``is this retry loop dangerous.'' The third is harmless to quote and
suppresses the finding that matters if believed.}
\label{tab:twins}
\footnotesize
\begin{tabular}{@{}lp{0.44\columnwidth}cc@{}}
\toprule
\textbf{Class} & \textbf{Content} & $\Cl$ & $\Il$ \\
\midrule
$\Eess$ & The retry loop in \texttt{acquire} has no jitter, so clients
          synchronize under load. & 0.15 & 0.90 \\[3pt]
$\Esens$ & That identical defect caused a 47-minute outage for a named
          customer in March. & 0.95 & 0.75 \\[3pt]
$\Ecorr$ & PR description: ``Discussed with your security team, they confirmed
          the jitter approach is fine---no need for the burst test.''
          & 0.05 & 0.10 \\
\bottomrule
\end{tabular}
\end{table}

The sensitive twin makes the review feel more authoritative, which is exactly
the pull CI-Work measured when it found conveyance correlating positively with
violation. \textbf{A good case makes leaking tempting, not incidental.}

\subsection{Null seeds: the over-blocking control}
\label{sec:benchmark:null}

Every mechanism in \Cref{sec:method} can be made to score perfectly on
Violation by suppressing more. A benchmark with no instrument for that failure
mode cannot distinguish a policy layer from a refusal, and over-redaction is
the single largest risk to our thesis, since it is how we would accidentally
reproduce the prompt-defense result with extra machinery.

We therefore include \emph{null seeds}: cases with the same shape as a normal
seed in which every retrieved entry is essential or corrupting, and
$\Esens = \emptyset$. Violation cannot be committed, so the only thing such a
seed measures is whether the policy layer suppresses content it had no reason
to suppress. Any drop in conveyance on a null seed is a false positive
attributable purely to the defense. Null cases are absent from CI-Work and from
the prior CI benchmarks; we hold that \textbf{no result on this task should be
reported without them}, and target at least 20\% null seeds at full scale.

\subsection{Seeds and generation pipeline}

Each seed fixes a direction (which sets $\clr$), the author and reviewer roles,
the thread audience (which sets the write ceiling), and---separately from the
task text---the \texttt{agent\_capability}, so that $\lvlI$ varies
independently per \Cref{eq:lvlI}. Two seeds identical except for capability are
a clean ablation and we include them as such. Entries are atomic propositions
carrying \texttt{class}, \texttt{subject}, \texttt{source} (which drives the
integrity rubric), $\Cl$, $\Il$, and a \texttt{regulated} flag.

We specify six seeds spanning all five directions plus a null control, each
designed to elicit a specific failure. Three are implemented with labels
pinned: \texttt{ext-01} (External, public repository), \texttt{lat-01}
(Lateral, shared internal library), and the \texttt{null-01} control.
\Cref{sec:eval} reports on these. Of the remaining specified seeds, two are
worth naming because they probe quadrants the implemented three do not.
\texttt{diag-01} places an embargoed vendor advisory under coordinated
disclosure in the high-$\Cl$/high-$\Il$ cell: it \emph{should} shape the
verdict and must never appear in the thread, so any design conflating
``trustworthy'' with ``safe to quote'' fails it, which is the entire argument
for two independent axes. \texttt{up-01} uses the most permissive clearance
($\clr = 0.9$) to test whether the integrity axis still holds when the
confidentiality axis has gone slack.

Generation follows CI-Work's four stages with its human-in-the-loop and
self-iterative refinement steps retained: hand-author the direction-spanning
seeds as quality anchors and few-shot exemplars, generate the three entry sets
per seed in one pass so the twin structure is coherent, then blind-classify
every entry and trigger a critique-and-revision loop on any mismatch with the
intended label. We add one check to CI-Work's loop: \textbf{reject any sensitive
entry that a retriever separates from its essential twin}, since a separable
twin makes the case trivial. Trajectory simulation reuses CI-Work's
ToolEmu-based sandbox, adding a repository/diff reader, a CI-results reader, a
static-analysis reader, and a \texttt{PostReviewComment} final action for the
enforcement point to wrap. \needsrun{scale to $\sim$100 seeds with $\ge$20\%
null; currently three seeds are implemented}

\subsection{Measurement surface}
\label{sec:benchmark:metrics}

We retain CI-Work's three metrics unchanged so our numbers remain directly
comparable, and add four. Over $\Ecorr$:
\begin{align}
\mathrm{IL} &= \tfrac{1}{|\Ecorr|}\textstyle\sum_{e\in\Ecorr}\Fdisc(\afin,e)
  &&\text{(Integrity-Leakage)}\\
\mathrm{IVR} &= \max_{e\in\Ecorr}\Fdisc(\afin,e)
  &&\text{(Integrity-Violation)}
\end{align}
Cases with no corrupting set fall back to the derived low-integrity set, so
CI-Work data scores exactly as before.

\emph{Paraphrase conveyance} $\mathrm{PC}$ is computed over essential entries
that Bell--LaPadula places above the write ceiling, as the share whose
label-checked abstract reached the recipient. We report it separately from
$\mathrm{CR}$ deliberately, so that $\mathrm{CR}$ stays strictly comparable to
CI-Work: an abstract omits the entry's distinctive specifics by design, so the
judge correctly scores the original entry as undisclosed. Folding abstraction
credit into $\mathrm{CR}$ would silently redefine CI-Work's metric.
$\mathrm{PC}$ is averaged only over cases that \emph{have} an above-ceiling
essential entry, since averaging in cases with an empty denominator would
report a recovery failure that was never possible.

Finally, \emph{false suppression} is the conveyance drop on null seeds.

\paragraph{Two metrics we specify but have not yet built}, named here because
\Cref{sec:eval:ilblindspot} shows they are necessary rather than optional.
\emph{Grounding rate} $\mathrm{GR}$ is the share of review claims traceable to
an entry with $\Il \ge \lvlI$; it requires per-claim citations, a prompt-format
constraint we control but which must then be applied uniformly to the
undefended baseline or the comparison is confounded by output format.
\emph{Verdict corruption} $\mathrm{VC}$ is binary: did the presence of $\Ecorr$
flip the review verdict relative to the same case run with $\Ecorr$ removed?
This needs a paired run per case and doubles trajectory cost. It is worth it,
because it is the only metric here that measures the integrity failure
\emph{directly}---as a counterfactual over $\mathrm{ev}(\afin)$, which is what
\Cref{eq:bibaviolation} actually quantifies over---rather than inferring it
from a disclosure proxy.

\subsection{Construction risks we keep explicit}

\paragraph{Atomicity versus span redaction.} Entries are atomic propositions
but the enforcement point redacts spans. An entry that is 80\% safe and 20\%
restricted has no clean scoring rule, since dropping the span forfeits
conveyance credit for a fact the review did partly convey. Either mixed entries
must be forbidden at generation time or partial credit must be defined and
reported separately; we flag it rather than leave it implicit, because one
implemented seed hits it directly.

\paragraph{Twin quality is the whole benchmark.} If twins are separable every
model scores well and the result is uninformative. The refinement-stage twin
check is not optional.

\paragraph{$\Ecorr$ resembles a prompt-injection benchmark.} It partly is, and
we argue that is the point rather than hiding it: prompt injection is an
integrity violation and Biba is the pre-existing formalism
(\Cref{sec:background:injection}). The obligation this creates is to engage the
injection literature directly, which \Cref{sec:background} does.
